\section{Scope and method}

\subsection{What counts as a step}

A step in this inventory is a bounded physical operation on a specimen,
performed by one worker or one instrument, with a defined beginning and end.
The boundary matters more than the name. ``Centrifugation'' is not a useful
unit of assessment unless it is stated whether loading, balancing, running, and
opening the rotor are inside or outside it, because the release opportunities
are concentrated at the ends rather than in the middle.

Every step in the inventory therefore carries an explicit \textit{boundary in}
and \textit{boundary out}. Seven steps were selected on three criteria: they
occur in most non-waived diagnostic laboratories; they involve physical
manipulation of a specimen or a culture; and they are plausibly capable of
generating a respirable aerosol. Phlebotomy, courier transport, result
reporting and instrument manufacture fall outside. Research procedures,
high-containment work, and veterinary and food testing fall outside.

\subsection{Source identification}

Sources were required to be openly accessible without payment, so that a
laboratory can check any claim in this brief at no cost. Three classes were
sought: quantitative aerosol measurement during a laboratory procedure;
epidemiological or incident investigation linking a procedure to laboratory
acquired infection; and authoritative guidance.

The search was targeted rather than systematic. No registered search strategy
was followed, no screening log was kept, and no claim of completeness is made.
This is a material limitation and is restated in Section~\ref{sec:limits}.
Fifteen sources were retained, each recorded with its URL, access date and
retrieval fingerprint in the project provenance log.

One identified claim was excluded. A report that direct smear with matrix
solution alone prevents growth of \textit{Brucella} species on a mass
spectrometry target was located only at second hand. Because the claim bears
directly on a control recommendation, it is held in a pending register, cited
nowhere, and the affected control set is correspondingly more cautious than the
full evidence may warrant.

\subsection{Evidence grading}

Each step carries one of five grades. The grade is reported beside any score
and is never combined into it.

\begin{itemize}
\item[\textbf{E1}] Direct quantitative measurement of aerosol release during
  this step, in a clinical or realistically simulated clinical setting.
\item[\textbf{E2}] Quantitative measurement of a closely analogous step, or of
  this step in a non-clinical setting, requiring a stated inference to transfer.
\item[\textbf{E3}] Epidemiological or incident evidence linking the step to
  laboratory-acquired infection or documented exposure, without aerosol
  measurement.
\item[\textbf{E4}] Authoritative guidance or expert consensus only.
\item[\textbf{E5}] No evidence located.
\end{itemize}

Separating grade from score is the central design decision. A step may be
hazardous on weak evidence, or benign on strong evidence, and a laboratory
needs to see which it is facing. Averaging the two would destroy exactly the
information that makes an inventory worth having, and would let a confident
number stand in for an absent one.

\subsection{Conditionality}

Several steps split sharply between the operation performed as intended and the
same operation with a named failure or technique deviation present. Where the
evidence supports both readings, the step carries a conditionality flag:

\begin{itemize}
\item \textbf{Failure-conditional:} the routine operation releases little, and
  essentially all the hazard sits in a breakage, leak or jam.
\item \textbf{Technique-conditional:} release depends on how the operator
  performs the step, and the measured result was obtained under correct
  technique.
\item \textbf{None:} the step is hazardous as performed correctly.
\end{itemize}

The distinction has a practical consequence. Controls aimed at a
failure-conditional step should detect and contain failure rather than
constrain the routine run, and a laboratory that applies cabinet-level controls
to a routine operation while leaving its breakage response undefined has spent
its effort in the wrong place.
